% Pursuing Stacks (cotes 134-2 to 134-8): archive cote 134-3, site page 19 (PDF page 20), typescript page 19.
% typescript. Draft compared: ps2.tex:851-922.
% AI-transcribed by Claude Opus 5.5 (claude-opus-5-5, adaptive thinking, effort high), Message Batches API, batch ps1 (msgbatch_01Rv5hCaNTtUt4JHx1yqoHPL), 2026-11-12; not yet checked by eye.
\scanpage{134-3/19}
\mspage{\hand{19}}

S),  such that one-point subsets are simplices and a subset of a simplex
is a simplex.  This category,  via the functor   $(S,K)\mapsto K$ ,  is equi-%
valent to the full subcategory of (Ord),  whose objects are those ordered
sets $K$,  such that for every $x\in K$,  the set $K_{\leq x}$  be isomorphic to the
\margin{\hand{(18)}}ordered set of \xout{\unsure{finite}} non empty subsets of some finite set (or ``simplex'').
Here \xout{\illegible{1}} the question \xout{\illegible{1}} whether  this category is a model category
in the technical sense \inserted{[of Quillen]} does'nt really arise\hand{(18)}\hand{,} because this category
does'nt \inserted{even} admit finite products -  rien à faire !
\margin{\unsure{(17)} 5.3. indeed, the notion of a maquette is not \unsure{an} algebraic structure species!}

23.  It may be about time to get back to stacks,  still I can't help go-%
ing on pondering about algebraic structures as models for homotopy types.
If we have any algebraic structure species,  giving \inserted{rise to} a category $M$ of
set-theoretic realizations,  the basic question here does'nt seem so
much whether $M$ is a model category for a suitable choice of the three
sets of arrows (fibrations,  cofibrations,  weak equivalences),  but
rather how to define a natural functor
\[ (*) \qquad M \hand{$\longrightarrow$} (\mathrm{Hot}) \hand{,} \]
where (Hot) is the category of usual homotopy types,  and see whether
via this functor (Hot) can be interpreted as a category of fractions
\hand{(}or ``localization''\hand{)}of $M$  -  namely,  of course,  by the operation of making
invertible those arrows in $M$ which are transformed into isomorphisms
in (Hot).  In any case,  if we have such a natural functor,  the natural
thing to do is to call those arrows ``weak equivalences''.  If we want $M$
\margin{\hand{(9)}}to be a category of models,  various examples suggest that the natural
thing again is to take as cofibrations the monomorphisms,\hand{(9)}\hand{,} and then
(expecting that the model categories we are going to meet will be closed
model categories) to define fibrations by the Serre-Quillen lifting pro-%
perty with respect to cofibrations (=monomorphisms) which are weak equi-%
valences.  This being done,  it becomes meaningful to ask if indeed $M$
is a category of models.

Now the reflections of the beginning of today's notes \deleted{p.14-15} \inserted{section 20}
suggest a rather natural way for describing a functor $(*)$,  which makes
\margin{\hand{(19)}}sense in fact,  in principle,  for any category $M$,  namely:  endow $M$ with
it's canonical topology (unless a still more natural one appears at hand\hand{(19)}
\margin{5.3. the canonical topology is \uline{not} always suitable, see \S 24 below}-  I am not sure there is any better one in the present context),  and
assume that for every  $X\in \mathrm{Ob}\, M$ ,  the pro-homotopy type of the induced
site  $M_{/X}$  is essentially constant,  i.e. can be identified with an
object in (Hot) itself.  We then get the functor $(*)$ in an obvious way.
It then becomes meaningful to ask whether this functor is a localization
functor.

When $M$ is defined in terms of an algebraic structure species,  it
